\documentclass[a4paper,11pt]{article}
\usepackage[utf8]{inputenc}
\usepackage[french]{babel}
\usepackage{amsmath}
\usepackage[osf,sups]{Baskervaldx} % lining figures
\usepackage[bigdelims,cmintegrals,vvarbb,baskervaldx]{newtxmath} % math font
\usepackage{graphicx}
\usepackage{xcolor}
\usepackage{url}
\usepackage[hidelinks]{hyperref}

% hyperrref configuration: remove ugly boxes and colors, keep urls in blue
\hypersetup{colorlinks,linkcolor=black,citecolor=black,urlcolor={blue!80!black}}


% paragraphs
\setlength{\parindent}{0em}
\setlength{\parskip}{1em}

% margins
\addtolength{\hoffset}{-3em}
\addtolength{\textwidth}{6em}
\addtolength{\voffset}{-3em}
\addtolength{\textheight}{6em}

% macros
\def\ud{\mathrm{d}}

\begin{document}

\begin{center}\Large
stage L3 : plan de travail semaines 1,2,3
\end{center}

{\bf EXERCICES DE GÉOMÉTRIE }

{\bf Exercice 0. (trigonométrie)}
Un filtre de cafè se plie à plat avec un angle droit à son sommet.  Une fois
deplié il prend une forme conique.  Calculez l'angle du cône.

{\bf Exercice 1. (géodésiques)}
Caractérisez géométriquement les géodésiques d'un cône.
Démontrez que toutes les géodésiques qui ne passent pas per le vertex sont
similaires, et calculez le nombre d'intersections d'une géodésique en fonction
de l'angle du cône.

{\bf Exercice 2. (longueur Riemannienne)}
Une métrique locale sur $\R^d$ est une fonction lisse~$g:\R^d\to S_n^{+}$
(ensemble des matrices symétriques définies positives, assimilables à des
formes bilinéaires).  On écrit~$g_x(u,v)=u^T g(x)v$.  La longueur d'une
courbe paramétrée~$q:[a,b]\to\R^d$ est, par définition
\[
	\ell(q)=\int_a^b\sqrt{g_{q(t)}\left(\dot q(t),\dot q(t)\right)}\ud t
\]
Démontrez que la longueur est invariante par re-paramétrisation:
$\ell(q\circ h)=\ell(q)$ pour toute~$h:\R\to\R$ dérivable avec~$h'>0$.

{\bf Rappel.}
Soit~$X$ un espace vectoriel normé et~$f:X\to\R$ une fonction différentiable.
La dérivée directionnelle de~$f$ en~$x\in X$ dans la direction~$v\in X$ est,
quand elle existe,
\[
	D_vf(x) = \lim_{\epsilon\to 0}\frac{f(x+\epsilon v)-f(x)}\epsilon
\]
Si~$f$ a un extremum en~$x$, alors les dérivées directionnelles de~$f$ en~$x$
sont toujours~0.

{\bf Exercice 3. (action vs. longueur)}
Le plan euclidien est un cas particulier de l'exercice antérieur pour~$d=2$
et~$g=\left(\begin{smallmatrix}1&0\\0&1\end{smallmatrix}\right)$.  On considère
ainsi la longueur euclidienne~$\ell(q)=\int_a^b\left\|\dot q(t)\right\|\ud t$
et l'\emph{action}~$A(q)=\frac12\int_a^b\left\|\dot q(t)\right\|^2\ud t$.

{\bf 3.0.} Démontrez que les courbes qui minimisent l'action sont des droites
paramétrées à vitesse constante~($\dot q=c$).

{\bf 3.1.} Démontrez que si une courbe minimise l'action, alors elle minimise
aussi la longueur {\sl (indication: Cauchy-Schwarz)}.

{\bf 3.2.} Démontrez que les courbes qui minimisent la longueur sont des
droites (paramétrées à vitesse quelconque).

{\sl
Indication pour 3.0 et 3.2: utiliser des dérivées directionnelles de~$A$ et
de~$\ell$ sur des directions~$v$ qui sont des courbes qui s'annulent
sur~$t=a,b$, puis intégrez par parties.
}

{\bf Exercice 4. (équations d'Euler-Lagrange)}
Soit~$L:\R\times\R^d\times\R^d\to\R$ une fonction différentiable que l'on
appellera~\emph{le Lagrangien}.  L'énergie d'une courbe~$q:[a,b]\to\R^d$ est
\[
	E(q)=\int_a^b L\left(t,q(t),\dot q(t)\right)\ud t.
\]
Démontrez que si une courbe~$q$ est minimum local de l'énergie, alors elle
satisfait lesdites équations d'Euler-Lagrange
\[
	\frac{\partial L}{\partial q}-\frac{\ud}{\ud t}\frac{\partial
	L}{\partial\dot q} = 0
\]
(système de~$n$ équations ordinaires, une pour chaque coordonnée
de~$q$.)

{\bf Exercice 5. (symboles de Christoffel d'une surface)}
On considère une métrique locale sur le plan définie par un champ de matrices
symétriques définies positives
$(x,y)\mapsto\left(\begin{matrix}E(x,y)&F(x,y)\\F(x,y)&G(x,y)\end{matrix}\right)$.
Démontrez que les courbes de longueur minimale par rapport à cette métrique
(dites géodésiques) satisfont un système d'EDO de la forme
\[
\begin{cases}
\ddot x=\Gamma^x_{xx}\dot x^2+\Gamma^x_{xy}\dot x\dot y+\Gamma^x_{yy}\dot y^2\\
\ddot y=\Gamma^y_{xx}\dot x^2+\Gamma^y_{xy}\dot x\dot y+\Gamma^y_{yy}\dot y^2\\
\end{cases}
\]
avec six fonctions~$\Gamma^k_{ij}$, dites les symboles de Christoffel, que l'on
peut exprimer en termes des fonctions~$E,F,G$ et leurs dérivées premières.

\clearpage

{\bf EXERCICES DE PHYSIQUE }

{\bf Exercice 0. (systèmes conservatifs)}
Soit~$V$ un potentiel en~$R^d$.  Démontrez que si~$q$ est une trajectoire du
système (c.à.d. une solution de l'équation~$\ddot q=-\nabla V(q)$) alors
l'énérgie~$E=\tfrac12\left\|\dot q\right\|^2+V(q)$ est constante au long de la
trajectoire.

{\bf Exercice 1. (reparamétrisation d'une trajectoire)}
Soit~$t\mapsto q(t)$ une trajectoire d'un système de potentiel~$V$ et~$t=t(s)$
une réparamétrisation.  Écrivez l'équation satisfaite par la courbe~$s\mapsto
q(t(s))$ et interprétez-la géométriquement.

{\bf Exercice 2. (principe de Maupertuis)}
On fixe~$E\in\R$ et on définit~$S(q)=\sqrt{2(E-V(q))}$ sur l'enseble de~$\R^d$
determiné par la condition~$E\ge V$.  Demontrez que les courbes qui minimisent
l'action~$A[q]=\int S$ sont des reparamétrisations de trajectoires.

\end{document}


% vim:set tw=79 sw=2 ts=2:
